\begin{table}[htbp]
\centering
\caption{Standardized Effect Sizes}
\label{tab:sde}
\begin{tabular}{lcccccc}
\toprule
Outcome & $\hat{\beta}$ & SE & SD($Y$) & SDE & SE(SDE) & Classification \\
\midrule
IMR (pooled, 5 cutoffs) & -0.321 & 0.413 & 13.53 & -0.0237 & 0.0305 & Small negative \\
IMR (cutoff 15,000) & -0.406 & 0.616 & 13.53 & -0.0300 & 0.0456 & Small negative \\
IMR (cutoff 30,000) & 0.030 & 0.900 & 13.53 & 0.0022 & 0.0665 & Null \\
IMR (cutoff 50,000) & 0.219 & 1.036 & 13.53 & 0.0162 & 0.0766 & Small positive \\
IMR (cutoff 80,000) & -0.082 & 1.045 & 13.53 & -0.0061 & 0.0772 & Small negative \\
IMR (cutoff 120,000) & -0.142 & 1.021 & 13.53 & -0.0105 & 0.0755 & Small negative \\
\bottomrule
\end{tabular}
\begin{tablenotes}[flushleft]\small
\item \textit{Notes:} \textbf{Country:} Brazil. \textbf{Research question:} Does crossing a constitutional population threshold that increases minimum city council size by two seats causally reduce infant mortality? \textbf{Policy mechanism:} Brazilian Constitution Article 29 (amended by EC 58/2009) mandates minimum numbers of city council members (vereadores) at discrete population thresholds; crossing a threshold mechanically adds two legislative seats, potentially increasing oversight of municipal health spending and primary care delivery. \textbf{Outcome definition:} Infant mortality rate---deaths of children under age one per 1,000 live births---from the Brazilian Mortality Information System (SIM) and Live Birth Information System (SINASC). \textbf{Treatment:} Binary indicator for municipal population exceeding the nearest constitutional threshold, which triggers a two-seat increase in minimum council size. \textbf{Data:} IBGE municipal population estimates, DATASUS SIM and SINASC vital statistics, 85,979 municipality-year observations across 5,401 municipalities, 2001--2020. \textbf{Method:} Multi-cutoff regression discontinuity design following Cattaneo et al.\ (2016); local polynomial estimation with triangular kernel and CCT (2014) optimal bandwidth; standard errors clustered at the municipality level. \textbf{Sample:} Restricted to five constitutional thresholds (15,000; 30,000; 50,000; 80,000; 120,000 inhabitants) where sufficient mass exists on both sides of the cutoff; municipality-years with zero live births excluded. SDE $= \hat{\beta} / \text{SD}(Y)$ where SD($Y$) is the pooled standard deviation of the infant mortality rate. Classification refers to magnitude, not statistical significance: Large ($|$SDE$| > 0.15$), Moderate ($0.05$--$0.15$), Small ($0.005$--$0.05$), Null ($< 0.005$).
\end{tablenotes}
\end{table}

